\subsection{Three further unit-terminal ordinary-eighth families}
\label{proof:Q5-C098}
Let the residual list be
\(\Sigma=(1,b_1,c_1,\ldots,b_k,c_k)\), with
\(1<b_1<c_1<\cdots<b_k<c_k\). The receiving root has an
outward unit arm. Every connector edge points from the terminal
to that root, so for any positive integer \(L\),
\begin{equation}\tag{Q98.0.1}\label{eq:Q98:0.1}
 H=E_rpT^L(E_tF).
\end{equation}
Here \(E_t,E_r\) are any genuine equal-pair products at their
actual roots. The following choices depend only on the residual
lengths:
\[
\begin{array}{c|c|c}
\Sigma&F&\E H^8\text{ upper}\\\hline
(1,3,c),\ c\ge5&qf_3h_c&1-d/601\\
(1,2,c),\ c\ge5&qf_2h_c&1-d/12000\\
(1,2,c,b,e),\ 4\le c<b<e&qf_2h_ch_bf_e&e^{-d/2}\\
(1,2,3,b,e),\ 5\le b<e&qf_2h_3h_bf_e&e^{-d/2}\\
(1,2,3,4,e),\ e\ge5&qf_2h_eh_3f_4&e^{-8d}.
\end{array}
\]
The crossed last row is used for every \(L\), including one.
The first two rows retain the stronger continuation estimates
\[
\begin{array}{c|c|c|c|c|c}
\Sigma&d\text{-domain}&L=2&L=3&L=4&L\ge5\\\hline
(1,3,c),\ c\ge5 &[0,1/8]&e^{-4d}&e^{-7d}&e^{-2d}&e^{-3d/4}\\
(1,2,c),\ c\ge5 &[0,1/5]&e^{-2d/5}&e^{-d/20}&e^{-d}&e^{-d}.
\end{array}
\]
For \(d\ge1/8,L\ge2\), the first row has the same-direction
bound \(e^{-3d/32}\). For \(d\ge1/5\), the second row compares
to the actual ordinary-eighth \(qf_2\) input at every \(L\).

\begin{proof}
Sections~\ref{part:unit-p:three-single}
and~\ref{part:unit-p:three-tail} prove the first row with its
complete signed means and all finite short vectors.
Section~\ref{part:unit-p:124-and-retained}, using the actual
\(9/4\)-retained packet, proves the second row and its strong
continuation rates. In these two rows the unit-connector inputs
are Sections~\ref{proof:Q5-C065} and~\ref{proof:Q5-C063}.

Section~\ref{part:unit-p:five-arms} proves the last three rows.
The first two one-step directions in that paragraph are the
actual multipair input of Section~\ref{proof:Q5-C067}; its
crossed one-step direction is proved directly. The high-domain
regrouping uses the shortest-two square family, already proved
in Section~\ref{part:unit-p:shortest-two}, with its original
terminal \(q\) and root \(p\).

Every comparison uses the same finite host and its original
measure. Restore decorations and full reversal by
Section~\ref{part:unit-p:decorations}; at \(d=0\) original
leaf integration gives equality.
\end{proof}
